\chapter{Conics}\label{geometry:ch:conics}

\begin{center}\itshape
It is written in the language of mathematics, and its characters are triangles, circles, and other geometric figures.
\upshape(Galileo Galilei, \emph{The Assayer}, 1623, translated by Stillman Drake)
\end{center}

Around 200 BC, Apollonius of Perga wrote eight books on the curves you get by cutting a cone with a plane. He gave them the names we
still use: ellipse, parabola and hyperbola, from Greek words for falling short, being equal and exceeding. Four of the books survive in Greek,
three more only in an Arabic translation made in Baghdad, and the eighth is lost. For eighteen centuries the conics were a beautiful and
useless subject. Then, within a few decades, they turned out to be the shapes of the world. Johannes Kepler, who had studied the special points of
the conics in his book on optics of 1604 and called them \emph{foci}, ``hearths'', after the burning point of a mirror, found in 1609 that Mars moves
on an ellipse with the Sun at a focus. Galileo showed that a thrown stone flies along a parabola. And in February 1640 Blaise Pascal, sixteen years old,
had a short \emph{Essay pour les coniques} printed in a few copies. Its first lemma contains the theorem that he called, according to later
tradition, the \emph{mystic hexagram}: the three pairs of opposite sides of a hexagon inscribed in a conic meet in three collinear points.

Projective geometry explains why the three curves of Apollonius are one: an ellipse, a parabola and a hyperbola are the same projective conic,
seen in charts where the line at infinity misses it, touches it or cuts it. This chapter proves that all nondegenerate conics with a point are
projectively the same, that five points determine a conic, Pascal's theorem with the help of the cross-ratio, and the theory of poles and polars
that turns a conic into a machine exchanging points and lines. We end with Poncelet's porism for two circles, a closing theorem as surprising
as Steiner's in Appendix~C, with a proof that uses a hidden invariant measure.

\begin{think}
Take two circles, one inside the other. From a point of the outer circle draw a tangent to the inner one, follow it to the outer circle, draw the
other tangent from there, and continue. Sometimes the zigzag closes after three steps, forming a triangle. If it closes for one starting point,
should it close for every starting point?
\end{think}

\section{Conics and quadratic forms}

Throughout the chapter $K$ is a field in which $1+1\ne0$, such as $\mathbb Q$, $\mathbb R$, $\mathbb C$ or $\mathbb F_p$ with $p$ odd.

\begin{definition}\label{geometry:def:conic}
Let $A$ be a nonzero symmetric $3\times3$ matrix over $K$, with quadratic form $q(x)=x^{\mathsf T}Ax$ and bilinear form $\beta(x,y)=x^{\mathsf T}Ay$.
The \emph{conic} $\mathcal C_A$ is the set of points $[x]\in\mathbb P^2(K)$ with $q(x)=0$. It is \emph{nondegenerate} if $\det A\ne0$.
\end{definition}

The condition $q(x)=0$ does not depend on the representative, since $q(\lambda x)=\lambda^2q(x)$. In the affine chart $z=1$ the equation becomes a
general equation of degree two, $ax^2+bxy+cy^2+dx+ey+f=0$. Degenerate conics are less interesting: for example $xy=0$ is a pair of lines and $x^2=0$ a
double line. Since $1+1\ne0$, the form $\beta$ is recovered from $q$ by $2\beta(x,y)=q(x+y)-q(x)-q(y)$.

\begin{proposition}\label{geometry:prop:conicline}
Let $\mathcal C$ be a nondegenerate conic.
\begin{enumerate}
\item $\mathcal C$ contains no line, and a line meets $\mathcal C$ in at most two points.
\item Let $P=[p]\in\mathcal C$. The line $T_P=\{[x]:\beta(p,x)=0\}$, the \emph{tangent} at $P$, meets $\mathcal C$ only at $P$. Every other line through
$P$ meets $\mathcal C$ in exactly one more point.
\end{enumerate}
\end{proposition}

\begin{proof}
(1) Suppose $\mathcal C$ contains the line given by a plane $V\subset K^3$. Then $q=0$ on $V$, so $\beta=0$ on $V\times V$ by the polarisation
formula, and $V\subset V^\perp=\{y:\beta(v,y)=0\text{ for all }v\in V\}$. Since $A$ is invertible, $V^\perp$ is the kernel of the map
$y\mapsto(\beta(v_1,y),\beta(v_2,y))$ for a basis $v_1,v_2$ of $V$, and these two linear forms $v_i^{\mathsf T}A$ are independent, so
$\dim V^\perp=1<2=\dim V$, a contradiction. Now let $L$ be any line, spanned by $u,v$. The restriction $q(su+tv)$ is a binary quadratic form in
$(s,t)$, not identically zero by what we just proved, so it has at most two roots $[s:t]$ in $\mathbb P^1(K)$.

(2) Note that $T_P$ is a line, because $Ap\ne0$, and it contains $P$, because $\beta(p,p)=q(p)=0$. Let $L$ be a line through $P$ and $R=[r]$ another
point of $L$. Then
\[
q(sp+tr)=2st\,\beta(p,r)+t^2q(r)=t\bigl(2s\,\beta(p,r)+t\,q(r)\bigr).
\]
The root $t=0$ is $P$. If $L=T_P$, then $\beta(p,r)=0$ and $q(r)\ne0$ (otherwise $q$ would vanish on $L$), so $P$ is the only point. If $L\ne T_P$,
then $\beta(p,r)\ne0$ and the second root $[s:t]=[q(r):-2\beta(p,r)]$ gives one more point, $[q(r)p-2\beta(p,r)r]\ne P$.
\end{proof}

How many points determine a conic? A symmetric matrix has six entries, so conics form a projective space of dimension $5$, and each point imposes one
linear condition: we expect five points to determine one conic.

\begin{theorem}\label{geometry:thm:fivepoints}
Through five points of $\mathbb P^2(K)$, no three of them collinear, there passes exactly one conic, and it is nondegenerate.
\end{theorem}

\begin{proof}
By Proposition~\ref{geometry:prop:frames} we may assume that four of the points are $[1:0:0]$, $[0:1:0]$, $[0:0:1]$ and $[1:1:1]$; call the fifth $[u:v:w]$. A
conic through the first three has no $x^2$, $y^2$, $z^2$ terms, so its equation is $ayz+bxz+cxy=0$. It passes through $[1:1:1]$ if and only if
$a+b+c=0$, and through the fifth point if and only if $avw+buw+cuv=0$. The fifth point is not on the line $z=0$ through the first two points, so
$w\ne0$, and in the same way $u,v\ne0$. The two linear conditions on $(a,b,c)$ are independent unless $(vw,uw,uv)$ is proportional to $(1,1,1)$, that
is unless $u=v=w$, which would make the fifth point equal to the fourth. So the solutions form a one-dimensional space: there is exactly one conic.

Its matrix has zero diagonal and off-diagonal entries $a/2$, $b/2$, $c/2$, and determinant $abc/4$. If $a=0$, then $b=-c\ne0$ and the second condition
reads $bu(w-v)=0$, so $v=w$; but then the fifth point lies on the line $y=z$ through $[1:0:0]$ and $[1:1:1]$, which is excluded. Similarly $b\ne0$
and $c\ne0$, so the conic is nondegenerate.
\end{proof}

\section{Parametrising a conic}

Proposition~\ref{geometry:prop:conicline} gives a bijection between the lines through a point $P$ of $\mathcal C$ and the points of $\mathcal C$: send a line to its
second intersection point, and the tangent to $P$ itself. Lines through $P$ form a projective line, so a conic with one point looks like
$\mathbb P^1(K)$. Here is the cleanest form of this.

\begin{theorem}[Normal form]\label{geometry:thm:normalform}
Every nondegenerate conic with at least one point is projectively equivalent to the conic $y^2=xz$, which is the image of
\[
\varphi:\mathbb P^1(K)\to\mathbb P^2(K),\qquad \varphi(t)=[1:t:t^2],\quad \varphi(\infty)=[0:0:1].
\]
\end{theorem}

\begin{proof}
Let $P\in\mathcal C$. A line through $P$ other than the tangent meets $\mathcal C$ in a second point $Q\ne P$. The tangents $T_P$ and $T_Q$ are distinct
lines (each meets $\mathcal C$ only at its own point), so they meet at a point $R$, and $R\notin\mathcal C$, since $R\in T_P\cap\mathcal C$ would force $R=P$
and similarly $R=Q$. The points $P,R,Q$ are not collinear (the line $PQ$ is not tangent at $P$, while $PR=T_P$), so we may take coordinates with
$P=[1:0:0]$, $R=[0:1:0]$, $Q=[0:0:1]$. Then $q(e_1)=q(e_3)=0$, and since $R$ lies on $T_P$ and $T_Q$, $\beta(e_1,e_2)=\beta(e_3,e_2)=0$. So the only
nonzero entries of $A$ are $\beta(e_1,e_3)=b$ and $q(e_2)=c$, and the equation is $cy^2+2bxz=0$, with $b,c\ne0$ because $\det A\ne0$. Replacing the
coordinate $x$ by $-\frac{c}{2b}x$ gives $c(y^2-xz)=0$. Finally, $y^2=xz$ with $x\ne0$ means $[x:y:z]=[1:t:t^2]$ for $t=y/x$, and with $x=0$ it means
$y=0$, the point $[0:0:1]$.
\end{proof}

Over $\mathbb R$, the circle $x^2+y^2=z^2$ is nondegenerate and has points, so \emph{every nondegenerate real conic with a point is a projective image
of the circle}. Over $\mathbb C$ every conic has points (a line meets it where a binary quadratic form vanishes, and that form has a root), so there is
only one nondegenerate complex conic. Over $\mathbb R$ the remaining case, a nondegenerate conic without points, is projectively $x^2+y^2+z^2=0$; this
follows from Sylvester's law of inertia in Appendix~F.

Over $\mathbb Q$ the theorem has an arithmetic meaning. If a conic with rational coefficients has one rational point, the lines with rational slope
through it give all the others, and the circle $x^2+y^2=1$ gives the Pythagorean triples (Exercise~\ref{geometry:ex:pythagorean}). But a conic may have no
rational point at all: $x^2+y^2=3$ is an example (Exercise~\ref{geometry:ex:nothree}). Deciding which conics have a rational point is a theorem of Legendre,
proved in Volume~12.

\section{Ellipse, parabola, hyperbola}

In the real affine plane, a nondegenerate conic with points meets the line at infinity in $0$, $1$ or $2$ points. It is called an \emph{ellipse}, a
\emph{parabola} or a \emph{hyperbola} accordingly. For $x^2/a^2+y^2/b^2=1$ the points at infinity solve $x^2/a^2+y^2/b^2=0$: none. For $y=x^2$,
homogenised to $yz=x^2$, the only point at infinity is $[0:1:0]$, where the line at infinity is tangent. For $xy=1$ they are $[1:0:0]$ and $[0:1:0]$,
the directions of the two asymptotes. So the three curves of Apollonius are one projective curve, seen in three different charts.

\begin{proposition}[Focus and directrix]\label{geometry:prop:focus}
Let $F$ be a point, $D$ a line not through $F$, and $e>0$. The set of points $P$ with $PF=e\cdot d(P,D)$ is a nondegenerate conic. It is an ellipse
if $e<1$, a parabola if $e=1$ and a hyperbola if $e>1$.
\end{proposition}

\begin{proof}
Take $F=(0,0)$ and $D$ the line $x=k$ with $k\ne0$. Squaring the condition gives $x^2+y^2=e^2(x-k)^2$, that is,
\[
(1-e^2)x^2+y^2+2e^2kx-e^2k^2=0,\qquad
A=\begin{pmatrix}1-e^2&0&e^2k\\0&1&0\\e^2k&0&-e^2k^2\end{pmatrix}.
\]
Expanding along the middle row, $\det A=-(1-e^2)e^2k^2-e^4k^2=-e^2k^2\ne0$. The conic has a point, the point of the segment from $F$ to $D$ dividing it
in the ratio $e:1$. Its points at infinity satisfy $(1-e^2)x^2+y^2=0$: there are none if $e<1$, one if $e=1$ and two if $e>1$.
\end{proof}

\begin{figure}[ht]
\centering
\fig{geometry-ch05-dandelin}
\caption{Dandelin's spheres, seen from the side. The cutting plane (blue) meets the cone in an ellipse; the two spheres inscribed in the cone and
tangent to the plane touch it at the foci $F_1$ and $F_2$.}
\label{geometry:fig:dandelin}
\end{figure}

\begin{beauty}
Why is the section of a cone by a plane an ellipse with two foci? In 1822 Germinal Pierre Dandelin gave a proof that needs no formula at all. Put in the
cone two spheres, each touching the cone along a circle and touching the cutting plane, one above and one below it (Figure~\ref{geometry:fig:dandelin}). Let
them touch the plane at $F_1$ and $F_2$. For a point $P$ of the section, the line from $P$ to the vertex of the cone touches the two spheres at
points $Q_1$ and $Q_2$ of the two circles of contact. Two tangent segments from a point to a sphere have the same length, so $PF_1=PQ_1$ and
$PF_2=PQ_2$, and therefore
\[
PF_1+PF_2=PQ_1+PQ_2=Q_1Q_2,
\]
the distance between the two circles measured along the cone, which is the same for every $P$. The points $F_1,F_2$ are the foci.
\end{beauty}

\section{The cross-ratio on a conic, and Pascal's theorem}

Four lines through a point $O$ have a cross-ratio: the cross-ratio of their intersections with any line $m$ not through $O$. It does not depend on
$m$, by Proposition~\ref{geometry:prop:perspectivity}: two such lines are related by central projection from $O$.

\begin{theorem}[Chasles and Steiner]\label{geometry:thm:chaslesconic}
Let $A,B,C,D$ be four distinct points of a nondegenerate conic $\mathcal C$. For a point $P\in\mathcal C$ different from them, the cross-ratio of the
lines $PA,PB,PC,PD$ does not depend on $P$.
\end{theorem}

\begin{proof}
By Theorem~\ref{geometry:thm:normalform} we may take $\mathcal C$ to be $y^2=xz$, with $A=\varphi(a)$, $B=\varphi(b)$, $C=\varphi(c)$, $D=\varphi(d)$ and
$P=\varphi(s)$. Assume first that $s\ne\infty$. For $t\ne s$,
\[
(1,s,s^2)\times(1,t,t^2)=(t-s)\,\bigl(st,\,-(s+t),\,1\bigr),
\]
so the line $P\varphi(t)$ is $[st:-(s+t):1]$. It meets the line $x=0$, which does not pass through $P=[1:s:s^2]$, at
$(st,-s-t,1)\times(1,0,0)=(0,1,s+t)$, the point with coordinate $s+t$ on that line. So the cross-ratio of the four lines is
$(s+a,s+b;s+c,s+d)=(a,b;c,d)$, since a translation is a M\"obius map (if one of $a,b,c,d$ is $\infty$, its line meets $x=0$ at $[0:0:1]$, which is
the point $\infty$ there, and the conclusion is the same). This does not depend on $s$. If $s=\infty$, the projective map $[x:y:z]\mapsto[z:y:x]$
preserves the conic and sends $\varphi(t)$ to $\varphi(1/t)$, which reduces to the first case and gives $(1/a,1/b;1/c,1/d)=(a,b;c,d)$.
\end{proof}

So four points on a conic have a cross-ratio of their own, and the parametrisation $\varphi$ is the dictionary: it equals the cross-ratio of the
parameters. The converse, due to Steiner, describes every conic as the locus of the intersections of corresponding lines of two pencils related
by a projective map; see Coxeter's \emph{Projective Geometry} in the sources. We need one more lemma.

\begin{lemma}\label{geometry:lem:perspectivity}
Let $L\ne M$ be lines meeting at $C$, and $h:L\to M$ a projective map with $h(C)=C$. Then $h$ is a central projection from some point $O$ on neither line.
\end{lemma}

\begin{proof}
Choose two more points $B,U$ of $L$, and let $V=h(B)$, $D=h(U)$; these are points of $M$ different from $C$. The lines $BV$ and $UD$ are distinct
(if they were equal, the line would meet $L$ in $B$ and $U$, so it would be $L$, but $V\notin L$). Let $O$ be their intersection. If $O$ were on $L$, it
would be $B$ (the only point of $BV$ on $L$) and also $U$; so $O\notin L$, and similarly $O\notin M$. The central projection from $O$ maps $C,B,U$ to
$C,V,D$, like $h$; both are projective maps $L\to M$, so they are equal by Proposition~\ref{geometry:prop:frames}.
\end{proof}

\begin{theorem}[Pascal]\label{geometry:thm:pascal}
Let $A,B,C,D,E,F$ be six distinct points of a nondegenerate conic. Then the points
\[
X=AB\cap DE,\qquad Y=BC\cap EF,\qquad Z=CD\cap FA
\]
are collinear.
\end{theorem}

\begin{proof}
By Theorem~\ref{geometry:thm:chaslesconic}, the pencils at $E$ and at $A$ through the points $B,C,D,F$ have the same cross-ratio:
$(EB,EC;ED,EF)=(AB,AC;AD,AF)$. Cut the first pencil by the line $BC$, which does not pass through $E$ (a line meets the conic at most twice). The four
lines meet it at $B$, $C$, $U=ED\cap BC$ and $Y$. Cut the second pencil by the line $CD$; the four lines meet it at $V=AB\cap CD$, $C$, $D$ and $Z$. So
\[
(B,C;U,Y)=(V,C;D,Z).
\]
The points $B,C,U$ are distinct ($U=B$ would put $B,D,E$ on a line), and so are $V,C,D$. Let $h:BC\to CD$ be the projective map with $h(B)=V$, $h(C)=C$,
$h(U)=D$. It preserves cross-ratios, so $(V,C;D,h(Y))=(B,C;U,Y)=(V,C;D,Z)$, and $h(Y)=Z$ (Theorem~\ref{geometry:thm:crossratio}). By
Lemma~\ref{geometry:lem:perspectivity}, $h$ is the central projection from $O=BV\cap UD$. But $BV$ is the line $AB$ and $UD$ is the line $DE$, so $O=X$. Hence
$Z$, the image of $Y$, lies on the line $XY$.
\end{proof}

\begin{figure}[ht]
\centering
\fig{geometry-ch05-pascal}
\caption{Pascal's hexagram. The hexagon $ABCDEF$ is inscribed in an ellipse; the meeting points $X,Y,Z$ of opposite sides lie on the Pascal line.}
\label{geometry:fig:pascal}
\end{figure}

Pascal's theorem has a dual. The tangent lines of a nondegenerate conic $\mathcal C_A$ form a conic in the dual plane: the tangent at $[p]$ is the
line $[Ap]$, and $\ell=Ap$ satisfies $p^{\mathsf T}Ap=0$ if and only if $\ell^{\mathsf T}A^{-1}\ell=0$. So the dual of Pascal's theorem is a theorem
about hexagons circumscribed about a conic: Brianchon's theorem of 1806 (Exercise~\ref{geometry:ex:brianchon}). When the conic degenerates into two lines,
Pascal's theorem becomes Pappus' theorem of Appendix~D (Exercise~\ref{geometry:ex:pascalpappus}).

\begin{remark}
Six points on a conic can be ordered into a hexagon in $60$ ways, which gives $60$ Pascal lines. They meet three by three in $20$ Steiner points,
and in other remarkable ways studied by Steiner, Kirkman, Cayley and Salmon; the configuration is described in Conway and Ryba's article listed in the
sources. A second proof of Pascal's theorem, by cubic curves, is in Appendix~G.
\end{remark}

\section{Poles and polars}

\begin{definition}
Let $\mathcal C=\mathcal C_A$ be nondegenerate. The \emph{polar} of a point $P=[p]$ is the line $[Ap]=\{[x]:\beta(p,x)=0\}$, and $P$ is the \emph{pole} of
this line.
\end{definition}

Since $A$ is invertible, this is a bijection between points and lines, and it preserves incidence in a twisted way: $P$ lies on the polar of $Q$
if and only if $\beta(q,p)=0$, if and only if $Q$ lies on the polar of $P$. This symmetry is \emph{La Hire's theorem}. The polar of a point of
$\mathcal C$ is its tangent.

\begin{proposition}\label{geometry:prop:polar}
Let $P\notin\mathcal C$.
\begin{enumerate}
\item A point $T$ of $\mathcal C$ lies on the polar of $P$ if and only if the tangent at $T$ passes through $P$. So if two tangents can be drawn from
$P$, the polar is the line through their points of contact.
\item If a line through $P$ meets $\mathcal C$ in two points $M,N$ and the polar of $P$ in $Q$, then $(M,N;P,Q)=-1$.
\end{enumerate}
\end{proposition}

\begin{proof}
(1) Both conditions say $\beta(t,p)=0$. (2) The point $P$ is not on its polar, since $\beta(p,p)=q(p)\ne0$; so the line meets the polar in one point
$Q=[r]$, with $\beta(p,r)=0$. Then $q(sp+tr)=s^2q(p)+t^2q(r)$, because the cross term vanishes. If $q(r)$ were $0$ the only root would be $s=0$,
giving one point; so $q(r)\ne0$ and the two roots are $t/s=\pm\lambda$ with $\lambda^2=-q(p)/q(r)$. In the coordinate $t/s$ on the line, $P=0$, $Q=\infty$,
$M=\lambda$, $N=-\lambda$, and $(\lambda,-\lambda;0,\infty)=\frac{0-\lambda}{0+\lambda}=-1$.
\end{proof}

So the polar of $P$ is the locus of the harmonic conjugates of $P$ with respect to the pairs of points where lines through $P$ cut the conic
(Figure~\ref{geometry:fig:polar}). Combined with the harmonic construction of Appendix~D, this lets you draw the tangents to a conic from an outside point
with a straightedge alone (Exercise~\ref{geometry:ex:tangentconstruct}).

\begin{figure}[ht]
\centering
\fig{geometry-ch05-polar}
\caption{The polar of $P$ (red) passes through the points of contact $T_1,T_2$ of the tangents from $P$. On every line through $P$ that cuts the
conic at $M$ and $N$, the polar cuts at the harmonic conjugate $Q$ of $P$.}
\label{geometry:fig:polar}
\end{figure}

\section{Poncelet's porism}

Let $\Gamma$ be a circle and $\omega$ a circle strictly inside it. For $x\in\Gamma$, of the two tangents from $x$ to $\omega$ take the one that has
$\omega$ on its right when you walk along it from $x$, and let $T(x)$ be the second point where it meets $\Gamma$. The sequence $x,T(x),T^2(x),\dots$
is a zigzag of chords of $\Gamma$ tangent to $\omega$. Poncelet stated in his \emph{Trait\'e} of 1822 that whether it closes up does not depend on the
starting point.

\begin{theorem}[Poncelet's porism for two circles]\label{geometry:thm:poncelet}
If $T^n(x)=x$ for one point $x\in\Gamma$, then $T^n(y)=y$ for every $y\in\Gamma$.
\end{theorem}

The proof finds a way of measuring arcs of $\Gamma$ for which $T$ is a rotation. For $x\in\Gamma$, let $t(x)>0$ be the length of the tangent segment
from $x$ to $\omega$; it is a continuous function, with $t(x)^2=|x-c|^2-r^2$ where $c$ and $r$ are the centre and the radius of $\omega$. Measure arcs
of $\Gamma$ by
\[
\mu(I)=\int_I\frac{ds}{t},
\]
where $ds$ is arc length.

\begin{lemma}\label{geometry:lem:poncelet}
$T$ is an orientation preserving homeomorphism of $\Gamma$ and $\mu(T(I))=\mu(I)$ for every arc $I$.
\end{lemma}

\begin{proof}
The tangent line depends continuously on $x$, so $T$ is continuous; it is bijective, since the inverse map is the same construction with the other
tangent. Let $x'$ be close to $x$, just after it counterclockwise, and let $M$ be the point where the chords $x\,T(x)$ and $x'\,T(x')$ cross. Two chords of
a circle cross exactly when their endpoints alternate around the circle, so $T(x')$ comes just after $T(x)$: $T$ preserves the orientation.

The triangles $Mxx'$ and $M\,T(x')\,T(x)$ are similar: their angles at $M$ are vertical angles, and the angles at $x$ and at $T(x')$ are inscribed
angles on the same arc from $x'$ to $T(x)$. Therefore
\[
\frac{|x'-x|}{|T(x')-T(x)|}=\frac{|M-x|}{|M-T(x')|}.
\]
Let $x'\to x$. The two chords are tangent to $\omega$ at points $P'$ and $P$, and $M$ is the meeting point of these tangents, at distance
$r\tan(\theta/2)$ from $P$, where $\theta$ is the angle between the radii $cP$ and $cP'$; so $M\to P$. Hence $|M-x|\to|P-x|=t(x)$ and
$|M-T(x')|\to|P-T(x)|=t(T(x))$. As chords and arcs are equivalent for nearby points, in an arc-length coordinate $T$ has the derivative
\[
T'=\frac{t\circ T}{t},
\]
a continuous positive function. By the change of variables formula,
$\mu(T(I))=\int_I\frac{T'(s)}{t(T(s))}\,ds=\int_I\frac{ds}{t(s)}=\mu(I)$.
\end{proof}

\begin{proof}[Proof of Theorem~\ref{geometry:thm:poncelet}]
Write $[x,y]$ for the counterclockwise arc from $x$ to $y$, and let $g(x)=\mu([x,T(x)])$. For $x'$ just after $x$, both $\mu([x,T(x')])=\mu([x,x'])+g(x')$
and $\mu([x,T(x')])=g(x)+\mu([T(x),T(x')])$ hold, and $\mu([T(x),T(x')])=\mu(T([x,x']))=\mu([x,x'])$ by the lemma and the orientation. So $g(x')=g(x)$:
the function $g$ is locally constant, hence constant on the connected circle, equal to some $g_0$. Fix $x_0\in\Gamma$ and let
$F(x)=\frac{2\pi}{\mu(\Gamma)}\mu([x_0,x])$, an angle defined modulo $2\pi$. Then $F$ is a homeomorphism from $\Gamma$ to the circle of angles, and
$F(T(x))=F(x)+\alpha$ modulo $2\pi$, with $\alpha=2\pi g_0/\mu(\Gamma)$. In the coordinate $F$, the map $T$ is the rotation by $\alpha$. A rotation
satisfies $T^n(x)=x$ for one point if and only if $n\alpha\in2\pi\mathbb Z$, and then it holds for every point.
\end{proof}

\begin{figure}[ht]
\centering
\fig{geometry-ch05-poncelet}
\caption{Poncelet triangles. The circles have radii $R=1$ and $r=0.36$, and the distance between their centres satisfies $d^2=R^2-2Rr$. Three
starting points give three triangles, each inscribed in the outer circle and circumscribed about the inner one.}
\label{geometry:fig:poncelet}
\end{figure}

When do the triangles close? In Appendix~A you proved Euler's formula $OI^2=R(R-2r)$ for a triangle with circumcentre $O$ and incentre $I$
(Exercise~\ref{geometry:ex:eulerOI}). Exercise~\ref{geometry:ex:chappleeuler} shows that the converse is the closing condition: the Poncelet zigzag between two
circles closes after three steps exactly when $d^2=R^2-2Rr$. For quadrilaterals the condition, found by Nicolaus Fuss, is
\[
\frac1{(R-d)^2}+\frac1{(R+d)^2}=\frac1{r^2},
\]
which a numerical calculation tests numerically.

\begin{remark}
Poncelet's theorem holds for any two nondegenerate conics in general position, over $\mathbb C$, and the proof is the same idea in disguise: the
pairs (point of the first conic, tangent line to the second through it) form an elliptic curve, a torus, and the Poncelet map is a translation on it.
Jacobi gave an analytic proof with elliptic functions in 1828, and Cayley found explicit closing conditions for every $n$; see Griffiths and Harris and
the book of Dragovi\'c and Radnovi\'c. The cubic curves of Appendix~H are these elliptic curves.
\end{remark}

\section{Exercises}

\begin{exercise}[$\star$]\label{geometry:ex:fivepoints}
Find the conic through $(0,0)$, $(1,0)$, $(0,1)$, $(1,2)$, $(2,1)$. What kind of conic is it, and what is its centre?
\end{exercise}

\begin{exercise}[$\star$]\label{geometry:ex:tangentellipse}
Show that the polar of $(x_0,y_0)$ with respect to the ellipse $x^2/a^2+y^2/b^2=1$ is the line $xx_0/a^2+yy_0/b^2=1$. Deduce the equation of the
tangent at a point of the ellipse.
\end{exercise}

\begin{exercise}[$\star$]\label{geometry:ex:parabolareflect}
Let $\mathcal P$ be the parabola with focus $F=(0,f)$ and directrix $y=-f$. Show that the tangent at $P\in\mathcal P$ is the perpendicular bisector of
$FH$, where $H$ is the foot of the perpendicular from $P$ to the directrix. Deduce that light arriving parallel to the axis is reflected through the focus.
\end{exercise}

\begin{exercise}[$\star$]\label{geometry:ex:pythagorean}
Find all rational points of $x^2+y^2=1$ by drawing lines through $(-1,0)$, and deduce the formula $(m^2-n^2,2mn,m^2+n^2)$ for the Pythagorean triples.
\end{exercise}

\begin{exercise}[$\star$]\label{geometry:ex:nothree}
Show that the conic $x^2+y^2=3$ has no rational point.
\end{exercise}

\begin{exercise}[$\star$]\label{geometry:ex:pascalpappus}
Let $\ell,\ell'$ be two lines, $A,B,C\in\ell$ and $A',B',C'\in\ell'$. Show that Pascal's statement for the hexagon $AB'CA'BC'$ inscribed in the
degenerate conic $\ell\cup\ell'$ is exactly Pappus' theorem.
\end{exercise}

\begin{exercise}[$\star$]\label{geometry:ex:brianchon}
State and prove Brianchon's theorem: if a hexagon is circumscribed about a nondegenerate conic (its six sides are tangent lines), its three main
diagonals are concurrent.
\end{exercise}

\begin{exercise}[$\star\star$]\label{geometry:ex:conesection}
Let $\Pi$ be a plane of $\mathbb R^3$ not through $0$, with equation $ax+by+cz=1$. Show that the section of the cone $x^2+y^2=z^2$ by $\Pi$ is an ellipse
if $|c|>\sqrt{a^2+b^2}$, a parabola if $|c|=\sqrt{a^2+b^2}$ and a hyperbola if $|c|<\sqrt{a^2+b^2}$. (Think of $\Pi$ as an affine chart of
$\mathbb P^2(\mathbb R)$.)
\end{exercise}

\begin{exercise}[$\star\star$]\label{geometry:ex:tangentconstruct}
A conic is drawn on paper and $P$ is a point outside it. With a straightedge only, construct the two tangents from $P$.
\end{exercise}

\begin{exercise}[$\star$]\label{geometry:ex:concentric}
Let $\Gamma$ and $\omega$ be concentric circles of radii $R>r$. Show that $T$ is a rotation, and that the Poncelet zigzag closes into a convex $n$-gon
if and only if $r=R\cos(\pi/n)$. When does it close at all?
\end{exercise}

\begin{exercise}[$\star\star$]\label{geometry:ex:chappleeuler}
Let $\Gamma,\omega$ have radii $R>r$ and centres at distance $d$, with $\omega$ inside $\Gamma$. Using Euler's formula and Theorem~\ref{geometry:thm:poncelet},
show that if $d^2=R^2-2Rr$, then every point of $\Gamma$ is a vertex of a triangle inscribed in $\Gamma$ and circumscribed about $\omega$, and that if
one such triangle exists then $d^2=R^2-2Rr$.
\end{exercise}

\begin{exercise}[$\star$]\label{geometry:ex:apolloniuscircle}
Let $A\ne B$ be points and $k>0$, $k\ne1$. Show that the set of points $P$ with $PA=k\cdot PB$ is a circle (the circle of Apollonius), and that $A$ and
$B$ are inverse points with respect to it.
\end{exercise}

\section{Hints}

\begin{hint}{geometry:ex:fivepoints}
Write $ax^2+bxy+cy^2+dx+ey+f=0$ and solve the five linear equations. The centre is where the gradient vanishes.
\end{hint}
\begin{hint}{geometry:ex:tangentellipse}
The matrix is $\mathrm{diag}(1/a^2,1/b^2,-1)$; the polar of a point of the conic is the tangent.
\end{hint}
\begin{hint}{geometry:ex:parabolareflect}
$PF=PH$, so $P$ is on the perpendicular bisector of $FH$. Check that this line meets the parabola only at $P$.
\end{hint}
\begin{hint}{geometry:ex:pythagorean}
The second intersection of $y=t(x+1)$ with the circle has rational coordinates when $t$ is rational, and conversely.
\end{hint}
\begin{hint}{geometry:ex:nothree}
Clear denominators and work modulo $3$.
\end{hint}
\begin{hint}{geometry:ex:pascalpappus}
List the three pairs of opposite sides of the hexagon.
\end{hint}
\begin{hint}{geometry:ex:brianchon}
The tangent lines form a conic in the dual plane; apply Pascal there.
\end{hint}
\begin{hint}{geometry:ex:conesection}
The line at infinity of the chart is the set of lines through $0$ parallel to $\Pi$. Count the lines of the cone in the plane $ax+by+cz=0$.
\end{hint}
\begin{hint}{geometry:ex:tangentconstruct}
Find two points of the polar of $P$ using two lines through $P$ and the harmonic construction of Appendix~D.
\end{hint}
\begin{hint}{geometry:ex:concentric}
The chord tangent to $\omega$ subtends a fixed angle at the centre.
\end{hint}
\begin{hint}{geometry:ex:chappleeuler}
For the first part you need one triangle. Look among isosceles triangles inscribed in a circle of radius $R$.
\end{hint}
\begin{hint}{geometry:ex:apolloniuscircle}
Use complex numbers and expand $|z-a|^2=k^2|z-b|^2$.
\end{hint}

\section{Solutions}

\begin{solution}{geometry:ex:fivepoints}
\idea Five points give five linear equations for six unknown coefficients.
\writeup The point $(0,0)$ gives $f=0$; $(1,0)$ gives $a+d=0$; $(0,1)$ gives $c+e=0$; then $(1,2)$ gives $a+2b+4c+d+2e=0$, i.e.\ $2b+2c=0$, and $(2,1)$
gives $4a+2b+c+2d+e=0$, i.e.\ $2a+2b=0$. So $a=c=-b$, $d=e=b$, and with $b=-1$ the conic is
\[
x^2-xy+y^2-x-y=0.
\]
It is symmetric in $x$ and $y$. Its points at infinity satisfy $x^2-xy+y^2=0$, which has no real solution (the discriminant is $1-4<0$), so it is an
ellipse. The gradient $(2x-y-1,\,2y-x-1)$ vanishes at $(1,1)$, the centre. As a check, $(2,2)$, the reflection of $(0,0)$ in the centre, is also on the curve.
\end{solution}

\begin{solution}{geometry:ex:tangentellipse}
\idea The polar is a matrix times a vector.
\writeup The conic is $q(x,y,z)=x^2/a^2+y^2/b^2-z^2$ with $A=\mathrm{diag}(1/a^2,1/b^2,-1)$. The polar of $[x_0:y_0:1]$ is $[x_0/a^2:y_0/b^2:-1]$, the line
$xx_0/a^2+yy_0/b^2=1$. For a point of the ellipse the polar is the tangent (Proposition~\ref{geometry:prop:conicline}), so the tangent at $(x_0,y_0)$ is
$xx_0/a^2+yy_0/b^2=1$.
\end{solution}

\begin{solution}{geometry:ex:parabolareflect}
\idea The tangent at $P$ is the line that the reflection exchanging $F$ and $H$ fixes.
\writeup The parabola is $x^2=4fy$. Let $P=(u,u^2/4f)$, so $H=(u,-f)$. The perpendicular bisector of $FH$ passes through the midpoint $(u/2,0)$ and is
perpendicular to $H-F=(u,-2f)$, so its equation is $u(x-u/2)-2fy=0$, that is $y=\frac{u}{2f}x-\frac{u^2}{4f}$. It passes through $P$, and
substituting into $x^2=4fy$ gives $x^2-2ux+u^2=(x-u)^2=0$: it meets the parabola only at $P$, so it is the tangent. The reflection in the tangent
exchanges $F$ and $H$, so it maps the line $PH$, which is vertical, to the line $PF$. A ray coming down parallel to the axis along $HP$ is reflected
along $PF$, into the focus. This is why satellite dishes and the mirrors of telescopes are parabolic.
\end{solution}

\begin{solution}{geometry:ex:pythagorean}
\idea Rational lines through a rational point cut the conic again at rational points.
\writeup Substituting $y=t(x+1)$ into $x^2+y^2=1$ gives $(x+1)\bigl((1+t^2)x-(1-t^2)\bigr)=0$, so the second point is
$\bigl(\frac{1-t^2}{1+t^2},\frac{2t}{1+t^2}\bigr)$. It is rational if $t$ is, and conversely $t=y/(x+1)$ is rational if the point is. So the rational points
are $(-1,0)$ and these. A Pythagorean triple $a^2+b^2=c^2$ gives the rational point $(a/c,b/c)$; with $t=n/m$ in lowest terms, the point is
$\bigl(\frac{m^2-n^2}{m^2+n^2},\frac{2mn}{m^2+n^2}\bigr)$, so $(a,b,c)$ is proportional to $(m^2-n^2,2mn,m^2+n^2)$. The precise statement for primitive
triples (with $m,n$ coprime of opposite parity) is proved in Volume~1, Chapter~9.
\end{solution}

\begin{solution}{geometry:ex:nothree}
\idea Squares modulo $3$ are $0$ and $1$, so a sum of two squares divisible by $3$ needs both divisible by $3$.
\writeup Suppose $x^2+y^2=3$ with $x,y$ rational. Writing $x=a/c$, $y=b/c$ with a common denominator gives integers with $a^2+b^2=3c^2$, and we may assume
$\gcd(a,b,c)=1$. Modulo $3$, $a^2+b^2\equiv0$, and since squares are $0$ or $1$ modulo $3$, this forces $a\equiv b\equiv0$. Then $9\mid a^2+b^2=3c^2$,
so $3\mid c^2$ and $3\mid c$, contradicting $\gcd(a,b,c)=1$.
\end{solution}

\begin{solution}{geometry:ex:pascalpappus}
\idea Read off the opposite sides.
\writeup The sides of the hexagon $AB'CA'BC'$ are $AB'$, $B'C$, $CA'$, $A'B$, $BC'$, $C'A$, and opposite sides are the first and fourth, the second and fifth,
the third and sixth. Their meeting points are $AB'\cap A'B=X$, $B'C\cap BC'=Z$ and $CA'\cap C'A=Y$, the three points of Pappus' theorem (Theorem~\ref{geometry:thm:pappusline}).
So Pascal's statement for this hexagon is Pappus' theorem. Our proof of Pascal does not apply to a degenerate conic, but Appendix~D proved this case
directly. Pascal and Pappus are two faces of one theorem about six points on a conic, as Appendix~G will show with cubic curves.
\end{solution}

\begin{solution}{geometry:ex:brianchon}
\idea Dualise Pascal, using that the tangent lines of a conic form a conic of the dual plane.
\writeup Let the sides $a,b,c,d,e,f$ of the hexagon be tangent to $\mathcal C_A$, with vertices $a\cap b$, $b\cap c$, \dots, $f\cap a$. As shown before the
statement, the tangent lines are the points of the nondegenerate conic $\mathcal C_{A^{-1}}$ of the dual plane. Apply Pascal's theorem there to the
six ``points'' $a,\dots,f$. The ``line'' joining $a$ and $b$ is the point $a\cap b$; the ``meeting point'' of the ``lines'' $ab$ and $de$ is the line
through the vertices $a\cap b$ and $d\cap e$, a main diagonal of the hexagon. Pascal says the three such ``points'' are ``collinear'', that is, the
three main diagonals pass through one point.
\end{solution}

\begin{solution}{geometry:ex:conesection}
\idea Apollonius' cone is the projective conic $x^2+y^2=z^2$, and a cutting plane is an affine chart.
\writeup Each point of $\Pi$ spans a line through $0$, and the lines through $0$ that meet $\Pi$ are exactly those not parallel to it. So $\Pi$ is an
affine chart of $\mathbb P^2(\mathbb R)$, whose line at infinity consists of the lines through $0$ in the plane $\Pi_0:ax+by+cz=0$. The section of the
cone is the conic $x^2+y^2=z^2$ seen in this chart, and its points at infinity are the lines of the cone lying in $\Pi_0$. Rotate about the $z$-axis so
that $b=0$, $a\ge0$; then $\sqrt{a^2+b^2}$ becomes $a$. If $c\ne0$, on $\Pi_0$ we have $z=-ax/c$, and the cone gives $y^2=(a^2/c^2-1)x^2$: only the
origin if $a<|c|$, one line ($y=0$) if $a=|c|$, two lines if $a>|c|$. If $c=0$, then $a>0$ and $\Pi_0$ is $x=0$, meeting the cone in the two lines $y=\pm z$.
So the section is an ellipse, a parabola or a hyperbola as stated: a plane less steep than the cone gives an ellipse, a plane parallel to a generator
a parabola, and a steeper plane a hyperbola.
\end{solution}

\begin{solution}{geometry:ex:tangentconstruct}
\idea Draw two secants from $P$; the four points they cut out form a quadrangle whose other two diagonal points lie on the polar of $P$.
\writeup Draw two lines through $P$ meeting the conic at $M_1,N_1$ and $M_2,N_2$. Let $R=M_1M_2\cap N_1N_2$ and $S=M_1N_2\cap N_1M_2$, and let the line $RS$ meet
$M_1N_1$ at $C$ and $M_2N_2$ at $D$. Apply Theorem~\ref{geometry:thm:harmonic} to the quadrangle with $E=M_1$, $F=M_2$, $G=N_1$, $H=N_2$: then $EF\cap GH=R$,
$EH\cap FG=S$, $EG\cap RS=C$ and $FH\cap RS=D$, so $(R,S;C,D)=-1$. Project the line $RS$ onto the line $M_1N_1$ from $M_2$: $R$ goes to $M_1$, $S$ to $N_1$,
$C$ to itself and $D$ to $M_2D\cap M_1N_1=P$. So $(M_1,N_1;C,P)=-1$, hence $(M_1,N_1;P,C)=-1$ (exchanging the last two points inverts the cross-ratio), and
$C$ lies on the polar of $P$ by Proposition~\ref{geometry:prop:polar}. In the same way $D$ does. Since $C\ne D$ (they lie on different lines through $P$ and are
not $P$), the polar is the line $CD=RS$. It meets the conic at the points of contact $T_1,T_2$, and the tangents are $PT_1$ and $PT_2$.
\end{solution}

\begin{solution}{geometry:ex:concentric}
\idea By symmetry every tangent chord looks the same from the centre.
\writeup A chord of $\Gamma$ tangent to $\omega$ is at distance $r$ from the centre, so it subtends the angle $2\theta$ at the centre, with $\cos\theta=r/R$.
All the chords are oriented the same way, so $T$ is the rotation by $2\theta$. The zigzag is a convex $n$-gon, going once around, when $2\theta=2\pi/n$,
that is $r=R\cos(\pi/n)$. It closes at all if and only if $\theta/\pi$ is rational: if $\theta=\pi k/n$ in lowest terms, it closes after $n$ steps,
going $k$ times around, as a star polygon.
\end{solution}

\begin{solution}{geometry:ex:chappleeuler}
\idea Euler's formula gives the necessary condition; for sufficiency, Poncelet reduces everything to finding a single triangle.
\writeup If a triangle is inscribed in $\Gamma$ and circumscribed about $\omega$, then $\Gamma$ is its circumcircle and $\omega$ its incircle, and Euler's
formula gives $d^2=OI^2=R^2-2Rr$. Conversely, assume $d^2=R^2-2Rr$. Consider isosceles triangles inscribed in a circle of radius $R$: their inradius
depends continuously on the apex angle, tends to $0$ when the triangle flattens, and equals $R/2$ for the equilateral triangle. Since $d^2\ge0$ we
have $r\le R/2$, so by the intermediate value theorem one of these triangles has inradius exactly $r$. By Euler's formula, the distance between its
circumcentre and incentre is $\sqrt{R^2-2Rr}=d$. Move it by an isometry taking its circumcentre to the centre of $\Gamma$ and its incentre to the
centre of $\omega$: it is now inscribed in $\Gamma$ and circumscribed about $\omega$. Its vertices give a point $x$ with $T^3(x)=x$, possibly after
reversing the order of the vertices to match the orientation of $T$. By Theorem~\ref{geometry:thm:poncelet}, $T^3(y)=y$ for every $y$. The zigzag from $y$
cannot close after one step, since a chord has two distinct ends, so it is a triangle.
\end{solution}

\begin{solution}{geometry:ex:apolloniuscircle}
\idea The condition is quadratic with the same coefficient for $x^2$ and $y^2$, and no $xy$ term: a circle.
\writeup With complex coordinates $z,a,b$, the condition $|z-a|^2=k^2|z-b|^2$ expands to
$(1-k^2)|z|^2-2\,\mathrm{Re}\bigl(z(\bar a-k^2\bar b)\bigr)+|a|^2-k^2|b|^2=0$. Dividing by $1-k^2$ and completing the square gives $|z-c|^2=\rho^2$ with
\[
c=\frac{a-k^2b}{1-k^2},\qquad \rho=\frac{k\,|a-b|}{|1-k^2|}.
\]
Then $c-a=\frac{k^2(a-b)}{1-k^2}$ and $c-b=\frac{a-b}{1-k^2}$ are positive multiples of the same vector, so $A$ and $B$ lie on the same ray from $c$, and
$|c-a|\,|c-b|=\frac{k^2|a-b|^2}{(1-k^2)^2}=\rho^2$: they are inverse in the circle (Definition~\ref{geometry:def:inversion}).
\end{solution}

\section*{Sources and further reading}

\begin{itemize}
\item Galileo Galilei, \emph{The Assayer} (1623), in \emph{Discoveries and Opinions of Galileo}, translated by S.~Drake, Doubleday, 1957, pp.~237--238
(the epigraph).
\item Apollonius of Perga, \emph{Conics}: Books I--IV in T.~L.~Heath, \emph{Apollonius of Perga: Treatise on Conic Sections}, Cambridge, 1896;
G.~J.~Toomer, \emph{Apollonius: Conics Books V to VII, the Arabic Translation of the Lost Greek Original}, Springer, 1990.
\item A.~Del Centina, \emph{On Kepler's system of conics in Astronomiae pars optica}, Archive for History of Exact Sciences 70 (2016), 567--589 (Kepler's
foci of 1604).
\item R.~Taton, article ``Pascal, Blaise'' in the \emph{Dictionary of Scientific Biography} (the \emph{Essay pour les coniques} of February 1640 and its
first lemma).
\item G.~P.~Dandelin, \emph{M\'emoire sur quelques propri\'et\'es remarquables de la focale parabolique}, Nouveaux m\'emoires de l'Acad\'emie royale de
Bruxelles 2 (1822), 171--200; K.~Kendig, \emph{Conics}, MAA, 2005.
\item H.~S.~M.~Coxeter, \emph{Projective Geometry}, 2nd ed., Springer, 1987 (Steiner's construction, Pascal, Brianchon, poles and polars);
M.~Berger, \emph{Geometry I and II}, Springer, 1987 (conics, Poncelet).
\item J.~Conway and A.~Ryba, \emph{The Pascal mysticum demystified}, Mathematical Intelligencer 34 (2012), 4--8.
\item J.-V.~Poncelet, \emph{Trait\'e des propri\'et\'es projectives des figures}, Paris, 1822; L.~D.~King, \emph{Three problems in search of a
measure}, American Mathematical Monthly 101 (1994), 609--628 (the invariant measure proof); V.~Dragovi\'c and M.~Radnovi\'c, \emph{Poncelet Porisms and
Beyond}, Birkh\"auser, 2011 (history, Chapple, Euler, Fuss, Jacobi, Cayley); P.~Griffiths and J.~Harris, \emph{On Cayley's explicit solution to
Poncelet's porism}, L'Enseignement Math\'ematique 24 (1978), 31--40.
\item A.~V.~Akopyan and A.~A.~Zaslavsky, \emph{Geometry of Conics}, AMS, 2007.
\end{itemize}
